\section{Complete ordinary proof of the signed-cap construction}
\label{sec:direct-signed-cap-full}

\subsection{Literal polynomial and logical scope}
The preferred polynomial is the fixed polynomial
\[
 P_{352}(z_0,\ldots,z_{351})
   =\sum_{r=0}^{218} R_r(z_0,\ldots,z_{351})^2
   \quad\text{in }\mathbb Z[z_0,\ldots,z_{351}].
\]
Here the ordered list of all 219 quadratic residuals $R_r$, including
its complete ordered variable list and every integer coefficient, is the
literal file \texttt{alternative/odd\_cap/equation.json}. A term
\texttt{[c,[i,j,...]]} means $c z_i z_j\cdots$, and an empty index list
means a constant. The complete collected coefficient list defining the
same polynomial is
\texttt{alternative/odd\_cap/expanded\_polynomial.json}. Its SHA-256 is
\begin{quote}\small\ttfamily
 a4b6e1f295015967ebf3cfe8f5bf397b1e1a9987b1d620fcd940f4cef8f1ce84.
\end{quote}
Thus the equation is an ordinary fixed polynomial equation, not a formula
with powers, factorials, Pell functions, or unbounded lists left as
primitive operations. Every unknown ranges over all of $\mathbb Z$.
The exact degree is four; there are 9,355 collected monomials and the
maximum absolute collected coefficient is 15,020,718. The simpler
355-variable sibling is defined analogously by the 220 residuals in
\texttt{alternative/signed\_cap/equation.json}. It has 9,370 monomials
and the same degree and coefficient height.

The proof below establishes the new arithmetic replacement in both
directions. Composed with the frozen parent's ordinary direct-arithmetic,
C222/refutation, and ZFC calibration, it gives
\[
 \mathrm{ZFC}\vdash
 \bigl(\exists z\in\mathbb Z^{352}\;P_{352}(z)=0\bigr)
 \longleftrightarrow
 \bigl(\text{there exists a finite ZFC refutation}\bigr).
\]
Consequently its no-zero sentence is equivalent in ZFC to
$\operatorname{Con}(\mathrm{ZFC})$. This is an ordinary mathematical
proof with the stated inherited calibration, not a newly completed
Lean or Rocq kernel proof or a newly printed object-language ZFC
derivation. Neither consistency nor inconsistency is assumed. The
construction does not simulate K215 or instantiate Jones's universal
polynomial. Variable-count comparisons here concern the direct-calculus,
relatively small-coefficient quartic architecture; other architectures
have much smaller variable counts and enormous fixed coefficients.

\subsection{The signed squared-cap theorem}
For $A\geq2$ define the integer Pell sequences by
\[
 X_0(A)=1,\qquad Y_0(A)=0,
\]
\[
 \begin{aligned}
 X_{j+1}(A)&=A X_j(A)+(A^2-1)Y_j(A),\\
 Y_{j+1}(A)&=X_j(A)+A Y_j(A).
 \end{aligned}
\]
Assume the family is nonempty, so $m\geq1$. Suppose
\[
 W\geq4,\qquad A-1>9W^3,\qquad
 1\leq k_i<W\ (1\leq i\leq m),\qquad m\leq W.
\]
Let $x_i,y_i$ be unrestricted integers satisfying
\[
 x_i^2-(A^2-1)y_i^2=1.
\]
Certify once that $V=Y_{2W}(A)$, and impose
\begin{align}
 V&=2W+1+s+\sum_{i=1}^m y_i^2,
       &&s\geq0, \label{eq:full-signed-cap}\\
 \sum_{i=1}^m(y_i-k_i)^2&=(A-1)Q,
       &&Q\in\mathbb Z. \label{eq:full-signed-pool}
\end{align}
Then
\[
 y_i=Y_{k_i}(A)>0,\qquad |x_i|=X_{k_i}(A)
 \quad(1\leq i\leq m).
\]
Conversely the canonical positive ordinates and either signs of their
abscissae extend to solutions of
\eqref{eq:full-signed-cap}--\eqref{eq:full-signed-pool}.

\paragraph{Soundness.}
Integer Pell classification gives indices $j_i\geq0$ and signs
$\epsilon_i\in\{1,-1\}$ such that
\[
 |x_i|=X_{j_i}(A),\qquad y_i=\epsilon_iY_{j_i}(A).
\]
Either sign may be used when $j_i=0$. Equation
\eqref{eq:full-signed-cap} gives $y_i^2<V$. Since $V>1$, it follows
that $|y_i|<V$, and strict monotonicity of $Y_j(A)$ gives
$j_i<2W$. In particular the endpoint $j_i=2W$ is excluded.

Modulo $A-1$, the recurrence gives $Y_j(A)\equiv j$. Therefore
\eqref{eq:full-signed-pool} implies
\[
 A-1\ \bigm|\ E:=\sum_{i=1}^m(\epsilon_i j_i-k_i)^2.
\]
Every absolute difference is less than $3W$, so
\[
 0\leq E<9mW^2\leq9W^3<A-1.
\]
Thus $E=0$. Each signed index $\epsilon_i j_i$ equals the positive
integer $k_i$, which forces $\epsilon_i=1$ and $j_i=k_i$.
Neither a negative ordinate nor a zero-index solution survives. This
argument deliberately does not orient the sign of $x_i$; that is a
separate consequence of the retained power rows below.

\paragraph{Constructive completeness.}
Set $y_i=Y_{k_i}(A)$. Every difference $y_i-k_i$ is divisible by
$A-1$, so the quotient $Q$ in \eqref{eq:full-signed-pool} is an
integer, indeed nonnegative. Put $T=Y_W(A)\geq1$. The norm and
positivity imply $X_W(A)>(A-1)T$. The double-index identity therefore
gives
\[
 V=Y_{2W}(A)=2X_W(A)T>2(A-1)T^2.
\]
Since each $k_i<W$, one has $y_i<T$ and
$\sum_i y_i^2\leq mT^2$. Consequently
\[
 V-\sum_i y_i^2>(2A-2-m)T^2>2W+1,
\]
using $A-1>9W^3$ and $m\leq W$. Hence the explicit restoration
\[
 s=V-2W-1-\sum_i y_i^2
\]
is nonnegative. This proves completeness without extrapolating from
finite examples or imposing an additional asymptotic assumption.

\subsection{A single ordinary indexed-Pell certificate supplies the cap}
Initially consider the simpler sibling. All letters below are integers,
with only $e$ expressly constrained to be nonnegative:
\begin{align}
 U^2-(A^2-1)V^2&=1, &
 H^2-(A^2-1)J^2&=1, \label{eq:full-cap-norms}\\
 S^2-(B^2-1)R^2&=1, &
 B&=1+4Vd, \label{eq:full-cap-comparison}\\
 J&=V^2(e+1), & e&\geq0, \label{eq:full-positive-aux}\\
 B&=A+Hf, & S&=U+Hg, \label{eq:full-cap-transfers}\\
 R&=2W+4Vh. && \label{eq:full-cap-index}
\end{align}
The aggregate cap gives $V\geq2W+1>0$. Thus $J\geq V^2>1$.
Auxiliary Pell classification implies $|H|\geq A\geq2$, in fact
$|H|\geq2A^2-1$.

No positivity guard on $B$ or $d$ is needed. If $d=0$, then $B=1$
and \eqref{eq:full-cap-transfers} would give $H\mid A-1$, contrary
to $|H|\geq A$. If $d\ne0$, then $B=1+4Vd$ gives
$|B|\geq4V-1\geq3$. Thus $|B|\geq2$. The coordinates
$d,B,f,S,R,g,h,H,U$ may all remain unrestricted.

The inherited signed indexed-Pell criterion applies with requested
index $2W$, principal ordinate $V$, auxiliary ordinate $J$, and
moduli $|H|$ and $4V$. Its hypotheses
\[
 1\leq2W\leq V,\qquad V^2\mid J,
\]
the three norm equations, and its four congruences are precisely the
rows above and the aggregate cap. The conclusion is
\[
 |U|=X_{2W}(A),\qquad V=Y_{2W}(A).
\]
For completeness, the signed extensions used in this application are
spelled out rather than silently assumed.

For negative $B$, define $X_j(B),Y_j(B)$ by the same polynomial
recurrence. For $j>0$,
\[
 X_j(B)=(-1)^jX_j(|B|),\qquad
 Y_j(B)=(-1)^{j-1}Y_j(|B|),
\]
while $X_0(B)=1$, $Y_0(B)=0$. Thus the classification at $|B|$
permits independently signed coordinates at the signed base $B$.
Polynomial base-congruence transfer and
$Y_j(B)\equiv j\pmod{B-1}$ continue to hold. The usual
residue-class proof consequently applies for $|B|\geq2$.
A signed principal $U$ may be replaced by $|U|$ while multiplying
$S$ by the same sign. Its norm and the congruence $S\equiv U\pmod H$
are preserved. The sign of $H$ changes no divisibility modulus.

Here is the precise short-interval step. Classification gives a
principal index $i>0$, an auxiliary index $n>0$, and a comparison
index $r\geq0$. Squared divisibility $V^2\mid J$ gives $V\mid n$,
and the size comparison gives $i\leq n$. The abscissa congruence,
with the independent sign of $S$ and the half-period sign-change
identity allowed, gives
\[
 r\equiv i\text{ or }-i\pmod{2n}.
\]
The ordinate congruence gives
\[
 r\equiv2W\text{ or }-2W\pmod{4V}.
\]
It follows that $i\equiv2W$ or $-2W\pmod{2V}$. Since
$1\leq i\leq V$ and $1\leq2W\leq V$, this forces $i=2W$,
including the endpoint case. This is the inherited indexed-Pell
argument applied to the stated signed domains.

For the simpler sibling's completeness, choose the canonical pair
at $2W$. The existing indexed-Pell completeness theorem supplies
$H,J,B,S,R$ and progression witnesses with $B>0$ and $e\geq0$.
They also satisfy the weaker unrestricted domains here. Only one
ordinary indexed certificate, not a stronger simultaneous theorem,
is required for this cap. An explicit strengthened construction for
the preferred sibling is provided below.

\subsection{The scale and seed premises are established before any powers}
The retained parent frame is the exact polynomial identity
\[
 W=2+\sum_{a\in\mathcal A}a^2,
\]
where $\mathcal A$ is the set of distinct actual bases, exponents,
and outputs of the nineteen retained power requests. The independent
checker verifies the fully expanded identity against the literal
41-atom inventory. Consequently every such atom has absolute value
less than $W$. Since $t=t_0+16$ with $t_0\geq0$, and $t$ is one
of these atoms, $W\geq258$. In particular the sixteen independent
seeds satisfy $m=16\leq W$.

The retained scale certificate, with $A,z\geq0$, is
\[
 A^2-\bigl((W+1)^2-1\bigr)\bigl(W(z+1)\bigr)^2=1.
\]
Its ordinate is positive. Classification at base $W+1$ and
$Y_\ell(W+1)\equiv\ell\pmod W$ imply a positive index $\ell$
divisible by $W$. Therefore
\[
 A\geq X_W(W+1)\geq(W+1)^W.
\]
For every $W\geq4$,
\[
 (W+1)^W-1>9W^3.
\]
At $W=4$ this is $624>576$. The ratio of successive exponential
bounds is greater than $W+2$, whereas the cubic ratio is at most
$(5/4)^3$, so the inequality persists by induction. It follows that
$A-1>9W^3$ without interpreting any requested power.

All sixteen seed exponents are already positive from retained
semantic guards and literal source rows. They are $t$, $K=2R$,
$\rho+t$, $F^2+t$, the nine positive encoded representatives,
$y+1$, $x+1$, and the positive derived-label expression. The
presemantic definitions are
\[
 \rho=F^2t^2,\qquad y=F^2+r_0+1,
 \qquad R=2+t+x+y+F^2,
 \qquad x,r_0,t_0\geq0.
\]
The sign of $F$ is not yet assumed. The derived-label expression
has a positive constant term and nonnegative constructor inputs.
The fresh literal lower-bound certificate proves every seed from
retained guards and the actual source rows; it uses neither a
removed seed-gap sign nor any conclusion of the cap theorem.
Thus the signed cap theorem applies before any factorial,
binomial, CRT range, or power output has been interpreted.

\subsection{Exact changed interface and quadratic implementation}
The frozen source is the square-anchored 489-integer direct
construction before its monic projection. For each of sixteen
independent seeds the principal norm and all power-congruence rows
are retained. The deleted rows are exactly its index-at-most-$Y$
row, comparison norm, comparison-abscissa congruence and
comparison-ordinate congruence; the old shared auxiliary norm,
shared comparison-base congruence and positive-product row are
also deleted. The sixteen scalar guards on $Y-k$ and two old
shared auxiliary scalar guards are removed. The precise original
row numbers and every deleted literal are recorded in
\texttt{cap\_certificate.json}.

No semantic sign, digit-pair guard, factorial ceiling, range row,
or Horn-flow source row is deleted. Unused arithmetic gates of
the old certificates remain temporarily as polynomial definitions,
then disappear under checked integral monic projection. They are
not premises of the new proof.

Every nonnegative slack is represented by the ternary form
\[
 \mathcal N(a,b,c)=a^2+b^2+c(c+1),\qquad a,b,c\in\mathbb Z.
\]
This takes only nonnegative values and represents every nonnegative
integer. Indeed Legendre's three-square theorem applies to $4n+1$;
exactly one of the three representing coordinates is odd, and a
permutation followed by division by four gives this form.

To remove any ambiguity about degree, introduce the four explicit
wires
\[
 a_V=AV,\qquad a_J=AJ,\qquad b_R=BR,\qquad v_2=V^2.
\]
For the preferred odd-quotient construction, the entire new cap
interface consists of the following quadratic residuals:
\[
\begin{aligned}
 &a_V-AV, &&a_J-AJ, &&b_R-BR, &&v_2-V^2,\\
 &U^2+V^2-a_V^2-1,
 &&H^2+J^2-a_J^2-1,
 &&S^2+R^2-b_R^2-1,\\
 &B-1-4Vd,
 &&J-v_2-2v_2e,
 &&B-A-Hf,\\
 &S-U-Hg,
 &&R-2W-4Vh,\\
 &V-2W-1-s-\sum_i y_i^2,\\
 &\sum_i(y_i-k_i)^2-AQ+Q,\\
 &s-a^2-b^2-c^2-c.
\end{aligned}
\]
Here $a,b,c$ are the three fresh witnesses for the aggregate
slack, not the comparison base $B$. Every $k_i$ is affine in
already existing source wires, so its squared contribution is
quadratic without extra summand variables. In the simpler sibling,
replace $J-v_2-2v_2e$ by $J-v_2-v_2e$ and add
$e-a_e^2-b_e^2-c_e^2-c_e$. All variable names, exact indices and
fresh guard witnesses are fixed in the supplied finite certificate.
The projection may subsequently remove some of these wires; their
restorations are supplied explicitly.

\subsection{All nineteen powers and their abscissa orientations}
After the cap theorem, each seed has
\[
 Y=Y_k(A)>0,\qquad X= X_k(A)\text{ or }-X_k(A).
\]
For a retained request at positive base $n$ and index $k$, the
unchanged row is
\[
 X+(n-A)Y=m+T_nu,\qquad T_n=2An-n^2-1.
\]
The frame supplies $1\leq n,k<W$ and $|m|<W$; its scale gives
$T_n>W^W+W$. If $X>0$, standard Pell substitution gives
$m\equiv n^k\pmod{T_n}$, and the strict interval bound yields
$m=n^k$.

If $X<0$, multiply the conjugate expressions and use the norm.
The result is
\[
 mn^k\equiv-1\pmod{T_n}.
\]
Since $|mn^k+1|<T_n$, this forces the integer identity $mn^k=-1$.
For $n\geq2$ it is impossible. At $n=1$ the sole inverse branch
is $m=-1$. Thus a request with base at least two orients its shared
principal pair. This reasoning uses no nonnegative quotient or
output assumption: all such quotients remain signed.

The complete noncircular recovery order is as follows.
\begin{enumerate}
\item The retained factorial remainder and slack give $D\geq1$.
The retained exact-binomial output slack gives $q=D+o+1\geq2$.
The request $q^t$ orients the $t$ pair. The request
$(\rho+1)^t$ sharing that pair therefore recovers its true value
$V_0\geq1$, even if $\rho+1$ initially equals one. The unchanged
ceiling identity $DF=V_0+\text{factorial slack}$ now gives
$F\geq1$, and consequently $V_0\geq2$.
\item The retained $B\geq0$ guard and
$C+F^2+1=B i_q$, with $C\geq1$, imply $B>0$. The request at
index $K$ for $B^K$ gives either $N=B^K>0$ or the sole inverse
case $B=1$, $N=-1$. In either case the second base
$N+V_0+1$ is at least two. That request orients the $K$ pair,
excludes the inverse branch, and recovers the common radix
$q=N(N+V_0+1)^K$.
\item Every other independent seed has a request at base $q$ or
$q+1$, both at least two. Those requests orient all remaining
principal pairs. The $t+1$ pair is obtained by the retained exact
Pell addition from the already oriented $t$ pair, not by an
unchecked extra seed.
\end{enumerate}
In this paragraph $B$ is the retained semantic binomial output;
it is distinct from the private comparison-base coordinate $B$
of the cap certificate, whose scope was local to that certificate.
The literal source assigns distinct variable indices to them.

For completeness, the entire request inventory in original output
labels is:
\begin{itemize}
\item \texttt{pow0}: $(\rho+1)^t$; \texttt{pow4}: $q^t$.
\item \texttt{pow1}: $B^K$; \texttt{pow2}: $(N+V_0+1)^K$.
\item \texttt{pow3}: $(q+1)^{\rho+t}$;
      \texttt{pow5}: $(q+1)^{F^2+t}$.
\item \texttt{pow6}, \texttt{pow8}, \texttt{pow9}, \texttt{pow10},
\texttt{pow11}, \texttt{pow13}, \texttt{pow14}, \texttt{pow16},
\texttt{pow17}: $(q+1)$ raised to the nine shifted encoded
representatives.
\item \texttt{pow12}: $q^{y+1}$;
      \texttt{pow15}: $q^{x+1}$.
\item \texttt{pow18}: $(q+1)$ raised to the positive
      derived-label expression.
\item \texttt{pow7}: $q^{t+1}$, using the derived addition pair.
\end{itemize}
This is nineteen requests at sixteen independent indices and one
derived index. The literal audit checks retention of every row,
its exact exponent binding, and coverage of every independent seed.

\subsection{Composition with the unchanged direct arithmetic}
Once the powers are recovered, the twelve paired guards of the
parent restore all digit remainder and slack signs from their
nonnegative sums $q^k-1$. The parent's ordered factorial proof
then gives $F=t!$, identifies the reciprocal binomial product,
and validates every CRT range digit. Its amplified reciprocal CRT
argument recovers each bounded local C222/Horn-flow row and the
derived label bound. The multiplicity-flow argument extracts a
finite refutation, to which the ordinary C222/ZFC calibration
applies. These retained arguments are mathematical dependencies,
not conclusions inferred from finite regression tests.

Conversely, start from a finite refutation and choose the semantic
data exactly as in the parent's constructive completeness:
the horizon, factorial, Horn accumulators, beta-CRT data,
reciprocal product, powers, representatives and digit witnesses.
These data determine $W$. The retained scale construction supplies
$A$. Choose every seed canonically at $A$, and derive the $t+1$
pair by addition. The old positive-power completeness supplies all
retained power quotients. The signed-cap completeness proved above
supplies the aggregate slack and pooled quotient, and the single
cap-index certificate supplies its private witnesses. Add the
required ternary slack representations and retain all semantic
guards. Assign unused old witness coordinates arbitrarily, then
compute their remaining gate definitions. Every new raw residual
vanishes.

This reconstruction may replace the private old Pell witnesses.
It claims existence-preserving equivalence, not same-tuple equality
of discarded coordinates. The Horn factorization is the parent's
exact polynomial identity, and every subsequent elimination has
an explicitly checked integral monic restoration. Thus the final
finite polynomial has exactly the stated existence status in
both directions.

\subsection{Odd auxiliary quotients: the preferred 352-variable form}
Replace only \eqref{eq:full-positive-aux} by
\[
 J=V^2(2e+1),\qquad e\in\mathbb Z.
\]
The single sign predicate for $e$ and its three witnesses are
removed. The remaining interface is unchanged. The actual emitted
polynomial, not merely the local count, has 352 variables and
219 quadratic residuals.

\paragraph{Soundness with a signed auxiliary ordinate.}
For integer $e$, the number $2e+1$ is nonzero and has absolute
value at least one. Hence $|J|\geq V^2>1$. The auxiliary norm
still gives $|H|\geq A$. Every use of $J$ in the signed
indexed-Pell theorem is through its absolute-value classification
or the divisibility $V^2\mid J$; its sign is not used. The
cap proof therefore remains valid with $|J|$ in place of $J$
in the lower bound, and no other premise changes.

\paragraph{Constructive odd quotient and the full private certificate.}
Let $r=2W$, and take the canonical cap pair
\[
 U=X_r(A),\qquad V=Y_r(A).
\]
The doubling identities make $V$ even and $U$ odd. Choose
\[
 H=X_{rV}(A),\qquad J=Y_{rV}(A)=VY_V(U),
\]
where the last equality is the Pell composition identity.
Polynomial expansion of $Y_V(U)$ gives an integer $T$ with
\[
 Y_V(U)=V U^{V-1}+(U^2-1)T.
\]
As $U^2-1=(A^2-1)V^2$, this yields
\[
 \frac{J}{V^2}=\frac{Y_V(U)}{V}
   =U^{V-1}+(A^2-1)VT.
\]
The right side is a positive odd integer: $U$ is odd and $V$
is even. Thus
\[
 e=\frac{J/V^2-1}{2}\in\mathbb Z
\]
provides the required auxiliary row.

Moreover $H=X_V(U)$ is odd, since $U$ is odd and $U^2-1$ is
even. The norm gives $\gcd(H,J)=1$; because $V\mid J$, it
follows that $\gcd(H,4V)=1$. The ordinary Chinese remainder
theorem therefore supplies arbitrarily large positive $B$ such
that
\[
 B\equiv A\pmod H,\qquad B\equiv1\pmod{4V}.
\]
Choose one and set $S=X_r(B)$, $R=Y_r(B)$. Polynomial
base transfer gives $S\equiv U\pmod H$, and base-one
reduction gives $R\equiv r\pmod{4V}$. The corresponding
integers $d,f,g,h$ satisfy every remaining cap row. This is an
explicit construction of the whole private certificate; no
unproved odd-quotient version of an inherited completeness
statement has been assumed.

The saving is precisely three coordinates and one quadratic
residual. Only the aggregate slack needs a ternary nonnegative
predicate. The simpler sibling remains independently supplied.

\subsection{Reproduction, finite checks, and attribution}
The complete local ordinary proof is
\texttt{alternative/PROOF.md}. To rebuild the preferred endpoint,
run \texttt{python build\_signed\_cap.py NEW --odd-auxiliary};
then run \texttt{python check\_signed\_cap.py NEW}.
Omitting the flag builds the simpler sibling. Both procedures
also pass with Python optimization enabled. The producer uses
only standard-library arithmetic and frozen reviewed polynomial
utilities; no remote state is modified.

The checkers verify the exact parent deletion boundary, strict new
row and variable inventories, absence of the omitted guard, all
sixteen presemantic index bounds, the 41-atom frame identity,
the nineteen request bindings, every integral monic restoration,
and the full collected quartic. The independent Horn audit checks
17 affine square roots, 11 product links, 18 literal rewrites,
and the absence of any escaped factor. Each endpoint has 229
checked monic restorations. All four fresh builds, comprising
two endpoints in normal and optimized modes, reproduce every
producer artifact byte for byte. Altered-input, certificate,
extra-row and coefficient controls are rejected. Small Pell
examples and concrete weakened-hypothesis counterexamples
supplement these checks; they do not decide consistency or
replace any unbounded mathematical argument.

The natural pooled-cap principle is credited to the incoming
\path{ConZFC_Shared_Pell_Bound_261_2026-10-24} packet.
The new step is the squared aggregate at cap index $2W$, which
recovers the signs and indices jointly and removes all individual
principal-ordinate sign and index tests. The signed comparison-base
extension, parent paired guards, square factorial anchor,
Horn factorization, direct numerical calculus and ordinary ZFC
calibration retain their original attribution and source pins.
No claim of a globally minimal equation or literature-wide
historical priority is made.
